%% notes-tex/028-gene-essentiality/sections/6-open.tex
%% [[experiments.028-gene-essentiality]]
%% https://github.com/Mjvolk3/torchcell/tree/main/notes-tex/028-gene-essentiality/sections/6-open.tex

\section{Open\secstatus{todo}}
\label{sec:open}

The held-out number is close to saturated: the round-12 ensemble's interval, 0.893
[0.84, 0.94], spans what 195 genes with 31 essential can resolve, and every lever on the
transformer side has been moved (width, depth, the preprocessor, the budget, the class
weight, six encodings). What remains is attribution and the module.

\begin{itemize}
  \item The pathway multi-hot carries the round-9 gain (round 14); which pathways, and
    whether a pathway-level flux statistic from the solve would substitute for the
    annotation, is not measured.
  \item The round-16 cell (the round-11 affine on the round-9 backbone) and the round-9
    cell are one cell to the resolution of 195 genes (rounds 16 and 17); a larger
    held-out set, such as the genome-wide labels scored on genes the released split
    never touched, is what would separate them.
  \item The MLP, the graph convolution and the transformer on the same three inputs are
    one number as fold ensembles (0.886, 0.888, 0.893). What separates architectures on
    this endpoint is not measurable on 195 genes; the genome-wide labels on genes the
    released split never touched are the larger held-out set that could.
  \item Round 13: the pFBA warm start on the stoichiometry-only layer. If its growth
    column reads at the solve's level at initialization and keeps it, the projection
    route carries the solve and the anchored constraint set was the loss; if not, the
    alternations are.
  \item The layer that solves is feasible on the GEM's box and, with the capacity
    ceilings calibrated, on the anchored one (rounds 15b to 15d). What remains is
    variance: the forced column reaches 0.83 to 0.84 on two of nine seeds and 0.5 to 0.7
    on the rest, with the absorbing infeasible state removed by slack rather than by step
    size. A seed-robust construction is untested: a barrier on the ceilings, the gated
    form with the ceilings binding (0.80 at scale 100), or the solve on more genes than
    the 629 the split trains on.
  \item The solve's flux vector read through a head rather than its growth column, which
    on the ladder carries 0.11 to 0.13 more AUROC than the growth ratio and would take
    the biomass objective out of the readout; and the growth LP with an objective other
    than biomass, since the layer that solves accepts any objective column.
  \item The layer that solves has the wrong shape for a training loop: the program runs
    on the CPU inside the forward pass, the GPU waits on it (0 to 13 percent
    utilization sampled during round 19, against a 1.6 hour cell where the alternating
    layer takes 19 minutes), and the bytes moved are not the cost, the synchronization
    is. Three constructions would remove it, in order of how much they change: a
    per-genotype basis warm start of the simplex through HiGHS's own interface, which
    keeps the exact solve and the exact duals; a batched first-order solve on the GPU;
    a smooth interior-point solve at a fixed barrier, which changes what the layer
    computes and belongs to its own rounds.
  \item The module's own hyperparameters have not been moved on this endpoint: the
    $k_{cat}$ tier (measured against predicted) and the thermodynamic source. Given
    round 15, the $k_{cat}$ tier is the one that sets whether the capacity box holds a
    growing flux at all.
\end{itemize}
